\section{Background}
\label{sec:background}

This section fixes terminology and names the four criterion families the library
implements. Readers familiar with network inference can skip to
Section~\ref{sec:architecture}.

\subsection{Networks}

A rooted phylogenetic network is a rooted directed acyclic graph whose leaves are
labelled by sampled taxa. Nodes of in-degree one are tree nodes and represent
ordinary divergence; nodes of in-degree two or more are reticulation nodes and
represent an event in which a lineage derives material from more than one parent.
Edges entering a reticulation carry inheritance probabilities $\gamma$ summing to
one, giving the proportion of loci descending through each parent.

A network with no reticulation nodes is a tree, which is why \phynetpy\ does not
give trees a separate type. Two derived notions recur. A \emph{displayed tree} is
obtained by retaining one incoming edge at each reticulation and suppressing the
rest; a network on $r$ reticulations displays up to $2^r$ trees, each with a
probability given by the product of the retained $\gamma$ values. A \emph{blob}
is a biconnected component, and contracting each blob to a node yields the
\emph{tree of blobs}, which captures the part of the history that is unambiguously
tree-like. Network \emph{level} is the largest number of reticulations in any
single blob, and bounds how entangled the reticulate structure is.

Branch lengths are measured on more than one scale. Sequence-based methods work
in expected substitutions per site; coalescent methods work in coalescent units,
related through the population-size parameter $\theta = 4N\mu$. The two are not
interchangeable, and a network carrying lengths on one scale will produce
plausible but wrong answers if read on the other.

\subsection{Criteria}

\paragraph{Parsimony.} Minimising deep coalescences (MDC) counts the extra
lineages implied by embedding observed gene trees in a candidate species history
and prefers the history requiring fewest \cite{Than2009MDC}. It needs only
gene-tree topologies, is cheap, and assumes no probabilistic model. Under
polyploidy the criterion is applied to a multiply-labelled tree, in which a
polyploid taxon appears once per subgenome \cite{Yan2022Allop}.

\paragraph{Likelihood.} Under the multispecies network coalescent, the
probability of a gene tree given a network is obtained by summing over the ways
lineages can be routed through reticulations and coalesce within branches
\cite{Yu2014ML}. This uses all the information in the data and is the most
expensive of the criteria, since the sum grows quickly with taxa and
reticulations.

\paragraph{Pseudo-likelihood.} Replacing the full likelihood with a product over
subsets of taxa -- in practice rooted triplets -- yields a composite objective
that is far cheaper and retains enough signal to rank topologies
\cite{Yu2015MPL}. The cost is that the product is not a normalised probability of
the data, which is why it cannot serve as a posterior target.

\paragraph{Bayesian inference.} Sampling networks from a posterior requires
proposals that change dimension, since adding or removing a reticulation changes
the parameter count; reversible-jump MCMC supplies the machinery
\cite{Green1995RJMCMC, Hastings1970}. Samplers exist for gene trees as data, for
sequence alignments with gene trees treated as latent variables
\cite{Wen2018MCMCSEQ}, and for biallelic markers, where the marker likelihood is
computed exactly without gene trees \cite{Bryant2012SNAPP, Zhu2018BiMarkers}.

The grouping above is the reason \phynetpy's interface has the shape it does.
Across these four families the data type, the assumed generative process and the
objective vary independently, so the natural parameterisation of the method set
is a product of three axes rather than a list of names.
